% Unit 4 review sheet -- 2 pages.
\documentclass{apchem-worksheet}

\apcunit{4}
\apcunittitle{Chemical Reactions}
\apcdoctitle{Review Sheet}
\apcsubtitle{Everything on the exam traces to this page \textbullet\ exam Mon Oct 26}

\begin{document}
\apctitleblock

\section*{\sffamily The five tools}

\renewcommand{\arraystretch}{1.45}
\noindent\begin{tabular}{@{}p{3.9cm}p{9.8cm}@{}}
\toprule
\textbf{Change classifier} & physical $=$ same substances, IMFs only;
  chemical $=$ new substances, bonds made/broken. Dissolving a salt is
  physical; dissolving \ce{HCl(g)} is chemical (it ionizes)\\
\textbf{Particulate rules} & conserve atoms AND charge; soluble salts drawn
  as separated ions; precipitate drawn as a cluster; excess reactant stays
  in the picture\\
\textbf{Net ionic recipe} & split strong electrolytes only (soluble salts,
  strong acids, strong bases); keep solids, liquids, gases, weak
  electrolytes whole; cancel spectators\\
\textbf{Reaction types} & precipitation (insoluble product), acid--base
  (proton transfer), redox (oxidation states change)\\
\textbf{Solution stoich} & volume $\times$ molarity $\to$ moles $\to$
  ratio $\to$ moles $\to$ mass or volume. Check limiting reactant using the
  coefficients, not raw moles\\
\bottomrule
\end{tabular}

\section*{\sffamily Solubility, the working set}

\begin{itemize}[leftmargin=*,itemsep=0.18em]
  \item \textbf{Always soluble:} \ce{NO3-}, alkali metal cations,
        \ce{NH4+}. Check these first --- they settle most questions.
  \item \textbf{Usually soluble:} \ce{Cl-}/\ce{Br-}/\ce{I-} except with
        \ce{Ag+}, \ce{Pb^2+}; \ce{SO4^2-} except with \ce{Ba^2+},
        \ce{Pb^2+}, \ce{Ca^2+}.
  \item \textbf{Usually insoluble:} \ce{CO3^2-}, \ce{PO4^3-}, \ce{S^2-},
        \ce{OH-} --- unless the cation is on the always-soluble list.
\end{itemize}

\section*{\sffamily Redox in four lines}

\begin{itemize}[leftmargin=*,itemsep=0.18em]
  \item Assign states: element 0; O is $-2$; H is $+1$; sum to the charge.
  \item \textbf{OIL RIG} --- Oxidation Is Loss, Reduction Is Gain.
  \item Oxidized $\to$ \emph{reducing} agent. Reduced $\to$
        \emph{oxidizing} agent. (The labels name what it does to the other
        species.)
  \item Free elements on both sides is a reliable redox tell.
\end{itemize}

\section*{\sffamily The traps, one line each}

\begin{itemize}[leftmargin=*,itemsep=0.22em]
  \item Splitting a weak acid, a weak base, or water in a net ionic
        equation.
  \item Drawing soluble salts as bonded pairs, or the precipitate as free
        ions.
  \item Losing atoms between ``before'' and ``after'' boxes, or deleting the
        excess reactant.
  \item Naming the oxidizing agent when you mean the reducing agent.
  \item Using a 1:1 titration ratio for a diprotic acid.
  \item Dividing excess-reactant moles by one solution's volume instead of
        the \emph{total} volume.
  \item Treating a colour change alone as proof of reaction.
\end{itemize}

\section*{\sffamily Self-check --- 10 questions, exam conditions}

\begin{enumerate}[leftmargin=*,itemsep=0.4em]
  \item Physical or chemical: \ce{NH3} gas dissolving and ionizing?
        \answerblank[3cm]{chemical}
  \item Net ionic for \ce{Ba^2+} with \ce{SO4^2-}:
        \answerblank[5.9cm]{\ce{Ba^2+ + SO4^2- -> BaSO4(s)}}
  \item Is \ce{Ag2CO3} soluble? \answerblank[2.5cm]{no}
  \item Reaction type for \ce{Mg + 2 HCl}:
        \answerblank[2.5cm]{redox}
  \item Oxidation state of Cr in \ce{CrO4^2-}:
        \answerblank[2cm]{$+6$}
  \item Reducing agent in \ce{Zn + Cu^2+ -> Zn^2+ + Cu}:
        \answerblank[2.5cm]{Zn}
  \item Moles in \SI{30.0}{\milli\liter} of \SI{0.500}{\Molar}:
        \answerblank[3cm]{\SI{0.0150}{\mole}}
  \item Net ionic for \ce{HF} with \ce{NaOH}:
        \answerblank[5.5cm]{\ce{HF + OH^- -> F^- + H2O}}
  \item Volume of \SI{0.200}{\Molar} \ce{NaOH} to neutralize
        \SI{25.0}{\milli\liter} of \SI{0.100}{\Molar} \ce{H2SO4}:
        \answerblank[3cm]{\SI{25.0}{\milli\liter}}
  \item In a particulate drawing, aqueous \ce{CaCl2} appears as:
        \answerblank[5.3cm]{1 \ce{Ca^2+} : 2 \ce{Cl-}, separated}
\end{enumerate}

\begin{apcnote}[How to study for this exam]
Redo WS 4's FRQ cold --- the exam's free-response is its twin in shape and
length. Then practice the chain \emph{predict products $\to$ formula
equation $\to$ net ionic $\to$ limiting reactant $\to$ mass} until it runs
without stopping. For every ``explain'' item, name the particles and what
happened to them before naming the property.
\end{apcnote}

\end{document}
